\subsection{GROUP AND REPRESENTATION ASSOCIATED TO A COXETER MATRIX}

We retain the notation of the preceding numbers. Let \(W = W(M)\) be the group defined by the family of generators \((g_{s})_{s \in S}\) and the relations \(^{1}\)

\[
\left(g _ {s} g _ {s ^ {\prime}}\right) ^ {m \left(s, s ^ {\prime}\right)} = 1, \quad \text { for } s, s ^ {\prime} \in S, m \left(s, s ^ {\prime}\right) \neq + \infty .
\]

PROPOSITION 3. There exists a unique homomorphism

\[
\sigma : \mathrm{W} \rightarrow \mathbf {G L} (\mathrm{E})
\]

such that \(\sigma(g_s) = \sigma_s\) for all \(s \in S\) . The elements of \(\sigma(W)\) preserve the bilinear form \(B_M\) .

To prove the existence and uniqueness of \(\sigma\) , it is enough to show that \((\sigma_s \sigma_{s'})^{m(s, s')} = 1\) if \(m(s, s') \neq +\infty\) . Now, if \(s = s'\) , this follows from the fact that \(\sigma_s\) is of order 2; if \(s \neq s'\) , it follows from what we proved in no. 2. Finally, since the reflections \(\sigma_s\) preserve \(B_M\) , so do the elements of \(\sigma(W)\) .

Remark. 1) We shall see in no. 4 that \(\sigma\) is injective; the group W can thus be identified with the subgroup of \(\mathbf{GL}(E)\) generated by the \(\sigma_s\) .

PROPOSITION 4. a) The map \(s \mapsto g_s\) from S to W is injective.

\begin{enumerate}
\def\labelenumi{\alph{enumi})}
\setcounter{enumi}{1}
\item
  Each of the \(g_{s}\) is of order 2.
\item
  If \(s, s' \in \mathbf{S}\) , \(g_s g_{s'}\) is of order \(m(s, s')\) .
\end{enumerate}

Assertion a) follows from the fact that the composite map

\[
s \mapsto g _ {s} \mapsto \sigma_ {s}
\]

from S to \(\mathbf{GL}(E)\) is injective.

For \(b)\) (resp. \(c\) ), we remark that the order of \(g_{s}\) (resp. the order of \(g_{s}g_{s'}\) ) is at most 2 (resp. at most \(m(s, s')\) ). Since we have seen in no. 2 that the order of \(\sigma_{s}\) (resp. of \(\sigma_{s}\sigma_{s'}\) ) is 2 (resp. \(m(s, s')\) ), we must have equality.

In view of \(a)\) , \(S\) can be identified with a subset of \(W\) by means of the map \(s \mapsto g_s\) .

COROLLARY. The pair (W, S) is a Coxeter system with matrix M.

This is simply the content of properties \(b)\) and \(c),\) together with the definition of \(W\) .

Remark. 2) We have thus shown that every Coxeter matrix corresponds to a Coxeter group.

